\documentclass[a4paper,10pt,fleqn]{article}
\usepackage[margin=1in]{geometry}
\usepackage{fuzz}

\begin{document}
\title{Vox WebSocket Receive-Ahead: A Z Specification}
\author{Daemon session ordering and cancellation}
\maketitle

\section{Scope}

One accepted socket dispatches one handler at a time. Only fetch allows
receive-ahead. Label frames with consecutive ordinals in ASGI receive
order; payloads and digest slots are outside the model. A receive can
remain outstanding across fetch completion, including one whose result
is already ready but has not yet been inspected by the session.

\section{State}

\begin{zed}
  Phase ::= idle | handling | fetching | closed
\end{zed}

\begin{zed}
  ReadState ::= noRead | oneRead | deferredError | readyDisconnect
\end{zed}

\begin{zed}
  WaitRead ::= waitNoDisconnect | waitDisconnect
\end{zed}

\begin{schema}{Session}
  phase : Phase \\
  reader : ReadState \\
  received, dispatched : \nat \\
  queued : \seq \nat
  \where
  dispatched \leq received \\
  \# queued \leq 32 \\
  received = dispatched + \# queued \\
  \forall i : \dom queued @ queued(i) = dispatched + i \\
  \# queued = 32 \implies reader = noRead \\
  reader = readyDisconnect \implies
    (phase = fetching \lor (phase = idle \land queued = \langle \rangle)) \\
  phase = closed \implies reader = noRead
\end{schema}

\begin{schema}{Init}
  Session'
  \where
  phase' = idle \\
  reader' = noRead \\
  received' = 0 \\
  dispatched' = 0 \\
  queued' = \langle \rangle
\end{schema}

\section{Transitions}

Start a read only with no outstanding read, either directly from idle
with no backlog or while fetching with room left. Reusing an existing
read issues no new read.

\begin{schema}{StartRead}
  \Delta Session
  \where
  reader = noRead \\
  (phase = idle \land queued = \langle \rangle) \lor
    (phase = fetching \land \# queued < 32) \\
  reader' = oneRead \\
  phase' = phase \\
  received' = received \\
  dispatched' = dispatched \\
  queued' = queued
\end{schema}

\begin{schema}{BufferFrame}
  \Delta Session
  \where
  phase = fetching \\
  reader = oneRead \\
  \# queued < 32 \\
  reader' = noRead \\
  phase' = phase \\
  received' = received + 1 \\
  dispatched' = dispatched \\
  queued' = queued \cat \langle received' \rangle
\end{schema}

Both normal and fetch handlers start via the same dispatch operation.
The empty-queue branch consumes the outstanding read directly; the
nonempty branch dispatches the FIFO head and leaves any read untouched.

\begin{schema}{Dispatch}
  \Delta Session \\
  kind? : Phase \\
  frame! : \nat
  \where
  phase = idle \\
  kind? \in \{handling, fetching\} \\
  phase' = kind? \\
  dispatched' = dispatched + 1 \\
  frame! = dispatched' \\
  (queued \neq \langle \rangle \land frame! = queued(1) \land
      reader' = reader \land received' = received) \\
  \quad \lor (queued = \langle \rangle \land reader = oneRead \land
      queued' = queued \land reader' = noRead \land
      received' = received + 1) \\
  \forall i : \dom queued' @ queued'(i) = queued(i + 1)
\end{schema}

\begin{schema}{FinishHandler}
  \Delta Session \\
  waitRead? : WaitRead
  \where
  phase \in \{handling, fetching\} \\
  phase = fetching \implies waitRead? = waitNoDisconnect \\
  reader \neq readyDisconnect \\
  phase' = idle \\
  reader' = reader \\
  received' = received \\
  dispatched' = dispatched \\
  queued' = queued
\end{schema}

A non-disconnect receive failure is retained during a fetch. It surfaces
only when the next ordered receive has drained the buffer.

\begin{schema}{DeferReceiveError}
  \Delta Session
  \where
  phase = fetching \\
  reader = oneRead \\
  reader' = deferredError \\
  phase' = phase \\
  received' = received \\
  dispatched' = dispatched \\
  queued' = queued
\end{schema}

A disconnect becomes ready in this model when it is visible to the
ordered receive or in the set of completed tasks returned to the fetch
watcher. When that set also contains the fetch task, record the disconnect
first: \texttt{FinishHandler} is then disabled, so buffered requests
cannot start before \texttt{Disconnect} closes the session.

\begin{schema}{ObserveDisconnectReady}
  \Delta Session \\
  waitRead? : WaitRead
  \where
  waitRead? = waitDisconnect \\
  reader = oneRead \\
  phase = fetching \lor (phase = idle \land queued = \langle \rangle) \\
  reader' = readyDisconnect \\
  phase' = phase \\
  received' = received \\
  dispatched' = dispatched \\
  queued' = queued
\end{schema}

\begin{schema}{CloseSession}
  \Delta Session
  \where
  phase \neq closed \\
  phase' = closed \\
  reader' = noRead \\
  received' = received \\
  dispatched' = dispatched \\
  queued' = queued
\end{schema}

\begin{schema}{Disconnect}
  CloseSession
  \where
  reader = readyDisconnect
\end{schema}

\begin{schema}{ReceiveFailure}
  CloseSession
  \where
  phase = idle \\
  queued = \langle \rangle \\
  reader \in \{oneRead, deferredError\}
\end{schema}

At saturation, no read is outstanding, so the model does not allow
\texttt{Disconnect} until there is a new read. Normal cancellation and
handler failures instead use \texttt{CloseSession}, which tears down any
outstanding read and permits no subsequent dispatch.

\end{document}
